% unit U2
\section{The Proof in Miniature}\label{sec:mini}
\mathtgt

This section develops the whole proof of capture prediction on a fragment of the \treeK target. The fragment has closures, lets and calls, and capture sets made of term variables and capture variables. It has no casts, no objects beyond a unit literal, no boxes and no capture members. It is small enough to state every definition and to run one program by hand.

The fragment is not mechanized on its own. Its running example comes from the \Lean examples \lean{C1} and \lean{C6} of \treeK. Their typings are decided by the kernel and their run is built by hand. Section~\ref{sec:mech-mini} maps every statement of this section to its Lean counterpart. Section~\ref{sec:problem-ideas} named two ideas, and this section shows the first, use sets as data. The second idea, that closed capture evidence over a typed store relates leaves, needs capture members, which Section~\ref{sec:target} adds.

\FloatBarrier
\subsection{Overview}\label{sec:mini-overview}

Before the definitions, here is the running example and what the section proves about it. Section~\ref{sec:mini-leaves} defines capture sets, types and leaves, and Section~\ref{sec:mini-typing} evidence, typing and use sets.

\paragraph{The platform.}
Two capture variables, $\kappa_1$ and $\kappa_2$, stand for two primitive capabilities, a log and a console (Section~\ref{sec:cc-types}). A capture variable stands for itself, as a type variable does in System~F. A program reaches each through a capability value, a closure that it calls. The capabilities $\mathit{log}$ and $\mathit{console}$ are identity closures on $\Unit$, each annotated with its capture variable:
\[
\begin{array}{@{}lcl@{}}
v_{\mathit{log}} & = & \lamc{\cset{\kappa_1}}{w}{\Unit}{w}{\ev{elem}\,\cset{w}\,\cset{\kappa_1, w}} \\
v_{\mathit{console}} & = & \lamc{\cset{\kappa_2}}{w}{\Unit}{w}{\ev{elem}\,\cset{w}\,\cset{\kappa_2, w}}
\end{array}
\]
The annotation $\cset{\kappa_1}$ of $v_{\mathit{log}}$ says that calling it may use $\kappa_1$. The evidence in the brackets shows that its body uses nothing beyond the annotation and the parameter.

\paragraph{A closure that returns a closure.}
The closure $\mathit{c1}$ takes a unit, calls $\mathit{log}$, and returns a closure that calls $\mathit{console}$:
\[
\begin{array}{@{}lcl@{}}
v_{\mathit{inner}} & = & \lamc{\cset{\mathit{console}}}{z}{\Unit}{\mathit{console}\,z}{\ev{elem}\,\cset{\mathit{console}, z}\,\cset{\mathit{console}, z}} \\
v_{\mathit{c1}} & = & \lambda^{\cset{\mathit{log}}}(y : \Unit).\,(\letc{r}{\mathit{log}\,y}{v_{\mathit{inner}}}{\cempty}{\ev{elem}\,\cempty\,\cempty}) \\
& & \quad \ann{\ev{elem}\,\cset{\mathit{log}, y}\,\cset{\mathit{log}, y}}
\end{array}
\]
Values are named $v_{\ldots}$, so that a value and the location that stores it never share a name. The annotation of $v_{\mathit{c1}}$ is $\cset{\mathit{log}}$, the term variable and not the capture variable $\kappa_1$ behind it. It does not mention $\mathit{console}$. This is the reading of Section~\ref{sec:cc-closures}: the capture set of a closure bounds what calling it may use.

\paragraph{Two programs.}
The store $\sigma_0$ holds the two capture binders, the two capabilities, the literal $\unitv$ at $\mathit{unit}$, and $\mathit{c1}$. Two programs run from it:
\[
\begin{array}{@{}lcl@{}}
\sigma_0 & = & \kappa_1,\ \kappa_2,\ \mathit{log} = v_{\mathit{log}},\ \mathit{console} = v_{\mathit{console}},\ \mathit{unit} = \unitv,\ \mathit{c1} = v_{\mathit{c1}} \\
\mathit{prog}_1 & = & \letc{\mathit{inner}}{\mathit{c1}\,\mathit{unit}}{\mathit{inner}\,\mathit{unit}}{\cset{\mathit{console}, \mathit{log}}}{f_1} \\
\mathit{prog}_2 & = & \letc{\mathit{inner}}{\mathit{c1}\,\mathit{unit}}{\mathit{unit}}{\cempty}{\ev{capvar}\,\mathit{unit}}
\end{array}
\]
Both call $\mathit{c1}$ and bind the closure it returns to $\mathit{inner}$. The first program calls $\mathit{inner}$. The second drops it and returns $\mathit{unit}$. Each let declares in its brackets the use set of its body, $\cset{\mathit{console}, \mathit{log}}$ and $\cempty$, with evidence that the uses of the body fit it. Example~\ref{ex:mini-typing} builds the evidence $f_1$.

\paragraph{The claims.}
The run of $\mathit{prog}_1$ reaches a state that calls $\mathit{console}$ (Table~\ref{tab:mini-run}). That is allowed, because the initial use set of $\mathit{prog}_1$ has the leaf $\kappa_2$ (Section~\ref{sec:mini-leaves}). The use set of $\mathit{prog}_2$ has the one leaf $\kappa_1$. Theorem~\ref{thm:mini-effect} says that no state reached from $\mathit{prog}_2$ inspects a variable with leaf $\kappa_2$. This holds although the run of $\mathit{prog}_2$ stores the closure that would call $\mathit{console}$.

\paragraph{Roadmap.}
Section~\ref{sec:mini-canon} proves the one theorem about capture evidence, Theorem~\ref{thm:mini-cap}: evidence for $C \le D$ puts the leaves of $C$ among those of $D$. Section~\ref{sec:mini-machine} adds stores and the machine, and the rest follows the textbook route of Section~\ref{sec:problem-textbook}. Its inversion of subcapturing at a call becomes Lemma~\ref{lem:mini-store}, which reads the annotation of a closure off the store. Preservation gains a use-set half, Lemma~\ref{lem:mini-step}, whose cases discharge the evidence that each redex carries (Section~\ref{sec:mini-step}). Section~\ref{sec:mini-run} carries that lemma along a run and derives effect safety.

\FloatBarrier
\subsection{Capture sets, types and leaves}\label{sec:mini-leaves}

\begin{figure}[tp]
\small
\color{tgtcol}
\[
\begin{array}{@{}lrcl@{}}
\text{Term variables} & & & x, y, z, r, w \\
\text{Capture variables} & & & \kappa \\
\text{Capture atoms} & c & ::= & x \mid \kappa \\
\text{Capture sets} & A, C, D, U & ::= & \cset{c_1, \ldots, c_n} \qquad n \ge 0 \text{, a list} \\
\text{Shapes} & S & ::= & \top \mid \pit[x]{T}{T'} \\
\text{Types} & T & ::= & \capt{S}{C} \\
\text{Capture evidence} & f, g & ::= & \ev{refl}\,C \mid \ev{trans}\,f\,g \mid \ev{elem}\,C\,D \mid \ev{union}\,f\,g \mid \ev{capvar}\,x \\
\text{Terms} & t, u & ::= & x \mid v \mid x\,y \mid \letc{x}{t}{u}{U}{f} \\
\text{Values} & v & ::= & \lamc{A}{x}{T}{t}{g} \mid \unitv \\
\text{Contexts} & \Gamma & ::= & \cdot \mid \Gamma, x : T \mid \Gamma, \kappa : \cstar \\
\text{Stores} & \sigma & ::= & \cdot \mid \sigma, x = v \mid \sigma, \kappa \\
\text{Continuations} & K & ::= & \cdot \mid K \triangleright \frlet{x}{u}{U}{f} \\
\text{States} & \mathit{st} & ::= & \mst{\sigma}{K}{t}
\end{array}
\]
\[
\begin{array}{@{}lcl@{}}
\Unit & \mathrel{:=} & \capt{\top}{\cempty} \\
\CapT{\kappa} & \mathrel{:=} & \capt{(\pit[w]{\Unit}{\Unit})}{\cset{\kappa}} \\
\cset{c_1, \ldots, c_m} \cup \cset{d_1, \ldots, d_n} & \mathrel{:=} & \cset{c_1, \ldots, c_m, d_1, \ldots, d_n} \\
C \subseteq D & \mathrel{:\Leftrightarrow} & c \in D \text{ for every } c \in C \\
\Gamma(x) = T & \mathrel{:\Leftrightarrow} & \Gamma = \Gamma', x : T, \Gamma'' \text{ for some } \Gamma' \text{ and } \Gamma'' \\
\sigma(x) = v & \mathrel{:\Leftrightarrow} & \sigma = \sigma', x = v, \sigma'' \text{ for some } \sigma' \text{ and } \sigma'' \\
\annot{\lamc{A}{x}{T}{t}{g}} & = & A \\
\annot{\unitv} & = & \cempty \\
\weak{X} & \mathrel{:=} & X \text{ read under one more binder, which } X \text{ does not mention} \\
X[z \mapsto y] & \mathrel{:=} & X \text{ with } y \text{ for every free } z \text{, in its capture sets and evidence too} \\
\sext{\sigma}{\sigma'} & \mathrel{:\Leftrightarrow} & \sigma' = \sigma, \sigma'' \text{ for some } \sigma'' \\[6pt]
\leaves{\Gamma}{\cempty} & = & \cempty \\
\leaves{\Gamma}{\cset{c} \cup C} & = & \leaves{\Gamma}{c} \cup \leaves{\Gamma}{C} \\
\leaves{\Gamma}{x} & = & \leaves{\Gamma'}{C} \qquad \text{if } \Gamma = \Gamma', x : \capt{S}{C}, \Gamma'' \\
\leaves{\Gamma}{\kappa} & = & \cset{\kappa} \qquad \text{if } \Gamma = \Gamma', \kappa : \cstar, \Gamma'' \\
\caple{\Gamma}{C}{D} & \mathrel{:\Leftrightarrow} & \leaves{\Gamma}{C} \subseteq \leaves{\Gamma}{D} \\
\capeqv{\Gamma}{C}{D} & \mathrel{:\Leftrightarrow} & \caple{\Gamma}{C}{D} \text{ and } \caple{\Gamma}{D}{C} \\[6pt]
\uses{x} & = & \cset{x} \\
\uses{v} & = & \cempty \\
\uses{x\,y} & = & \cset{x, y} \\
\uses{\letc{x}{t}{u}{U}{f}} & = & \uses{t} \cup U \\
\uses{\cdot} & = & \cempty \\
\uses{K \triangleright \frlet{x}{u}{U}{f}} & = & \uses{K} \cup U \\
\uses{\mst{\sigma}{K}{t}} & = & \uses{t} \cup \uses{K} \\
\insp{x\,y} & = & x \\
\insp{t} & & \text{undefined if } t \text{ is a variable, a value or a let} \\
\insp{\mst{\sigma}{K}{t}} & = & \insp{t}
\end{array}
\]
\color{black}
\caption{The fragment. Capture atoms are term variables and capture variables. Capture sets are lists, $C \cup D$ is concatenation and $C \subseteq D$ says that every atom of $C$ occurs in $D$. $\top$ is the empty object type of report~1, and $\unitv$ is the literal with no witnesses and no fields. The brackets $\ann{\ }$ hold data that the typing rules read and erasure removes: the use set a let declares, its avoidance evidence, and the closing evidence of a lambda. $\pit[x]{T}{T'}$ binds $x$ in $T'$. A lambda binds $x$ in $t$ and $g$, not in $A$. A let and a let frame bind $x$ in $u$ and $f$, not in $U$. A context binds each variable once, and so does a store. The leaves of a capture set are the capture variables it reaches through the types of its term variables. The first equation that applies is used.}
\label{fig:mini-syntax}
\end{figure}

Figure~\ref{fig:mini-syntax} gives the fragment. A capture set is a list of capture atoms, read as the finite set of its atoms. A context binds both kinds of variable. A term binder $x : T$ gives a variable its type. A capture binder $\kappa : \cstar$ introduces a capture variable that stands for itself and for nothing else. No term binds a capture variable. In a run, capture binders come from the store. Every judgment is well scoped. The type of a binder mentions only the binders before it, and a judgment in $\Gamma$ mentions only variables that $\Gamma$ binds. Leaves are taken only of sets whose atoms are bound.

A type $\capt{S}{C}$ is a shape $S$ with a capture set $C$. The shapes of the fragment are $\top$, the empty object type of \Rone, and the function shape $\pit[x]{T}{T'}$, which binds its parameter $x$ in the codomain $T'$. The domain and the codomain are types, so each carries its own capture set. $\Unit$ is $\top$ with the empty set, and its value $\unitv$ is the literal with no witnesses and no fields. $\CapT{\kappa}$ is the type of a capability, a function on $\Unit$ that may use $\kappa$.

\paragraph{A capturing type.}
The type of $\mathit{c1}$ is
\[
T_{\mathit{c1}} = \capt{(\pit[y]{\Unit}{\capt{(\pit[z]{\Unit}{\Unit})}{\cset{\mathit{console}}}})}{\cset{\mathit{log}}} .
\]
Read it from the outside. The outer set $\cset{\mathit{log}}$ bounds what calling $\mathit{c1}$ may use. The codomain is the type of the closure that the call returns, and its set $\cset{\mathit{console}}$ bounds what calling that closure may use. Both sets name term variables. Why does the annotation name $\mathit{log}$ rather than the capture variable $\kappa_1$? Because $\mathit{log}$ is the capability that $\mathit{c1}$ calls, and what it stands for can be read off its type, $\CapT{\kappa_1}$. The two choices mean the same, since $\cset{\mathit{log}}$ and $\cset{\kappa_1}$ have the same leaves, as the next paragraph defines. The contexts of the example are
\[
\Gamma_0 = \kappa_1 : \cstar,\ \kappa_2 : \cstar,\ \mathit{log} : \CapT{\kappa_1},\ \mathit{console} : \CapT{\kappa_2},\ \mathit{unit} : \Unit
\qquad\text{and}\qquad
\Gamma_1 = \Gamma_0,\ \mathit{c1} : T_{\mathit{c1}} .
\]
The value $v_{\mathit{c1}}$ is typed in $\Gamma_0$, and $\Gamma_1$ is the context of the store $\sigma_0$ (Section~\ref{sec:mini-machine}).

\paragraph{Leaves.}
What a capture set stands for in a context is its leaves, the capture variables it reaches when each term variable is replaced by the capture set of its type. Figure~\ref{fig:mini-syntax} defines them by four equations. The leaves of a set are those of its atoms. A capture variable is its own leaf. A term variable has the leaves of the capture set of its type, computed in the part of the context before its binder. The recursion is well founded, because the capture set of a binder mentions only binders before it. For the same reason the leaves of a set do not change when the context grows at its end.

As Section~\ref{sec:cc-types} noted, a term variable is never a leaf, and a variable whose type has the empty capture set has no leaves. So the atom $\mathit{unit}$ adds nothing to a set, and $\cset{x}$ in a set does not say that $x$ is a capability. It says that the set covers whatever the type of $x$ captures.

We write $\caple{\Gamma}{C}{D}$ when every leaf of $C$ is a leaf of $D$, and read it as ``$C$ is covered by $D$ in $\Gamma$'', the documentation's phrase (Section~\ref{sec:cc-types}). We write $\capeqv{\Gamma}{C}{D}$ when $C$ and $D$ have the same leaves. Leaves are lists like capture sets, but only which atoms occur matters to $\preccurlyeq$ and $\approx$. Examples list leaves without duplicates.

\begin{example}[Leaves in the running example]\label{ex:mini-leaves}
In $\Gamma_1$,
\[
\begin{array}{@{}lcl@{}}
\leaves{\Gamma_1}{\cset{\mathit{c1}}} & = & \leaves{\Gamma_0}{\cset{\mathit{log}}} \ = \ \leaves{\Gamma_0}{\cset{\kappa_1}} \ = \ \cset{\kappa_1} \\
\leaves{\Gamma_1}{\cset{\mathit{unit}}} & = & \leaves{\Gamma_0}{\cempty} \ = \ \cempty \\
\leaves{\Gamma_1}{\cset{\mathit{console}}} & = & \cset{\kappa_2} \\
\leaves{\Gamma_1}{\cset{\mathit{c1}, \mathit{unit}, \mathit{console}, \mathit{log}}} & = & \cset{\kappa_1, \kappa_2} \\
\leaves{\Gamma_1}{\cset{\mathit{c1}, \mathit{unit}}} & = & \cset{\kappa_1}
\end{array}
\]
The first equation goes through $\mathit{log}$, whose type $\CapT{\kappa_1}$ has the set $\cset{\kappa_1}$. The last two sets are the use sets of $\mathit{prog}_1$ and $\mathit{prog}_2$ (Example~\ref{ex:mini-typing}). So $\kappa_2$ is a leaf of the first and not of the second, as the overview claimed. The last equation is built by hand. The others are on paper.
\end{example}

\FloatBarrier
\subsection{Evidence, typing and use sets}\label{sec:mini-typing}

\begin{figure}[t]
\small
\color{tgtcol}
\[
\begin{array}{@{}l@{\qquad}l@{}}
\rn{C-Refl} & \inferrule{ }{\clej{\Gamma}{\ev{refl}\,C}{C}{C}} \\[6pt]
\rn{C-Trans} & \inferrule{\clej{\Gamma}{f}{C_1}{C_2} \\ \clej{\Gamma}{g}{C_2}{C_3}}{\clej{\Gamma}{\ev{trans}\,f\,g}{C_1}{C_3}} \\[6pt]
\rn{C-Elem} & \inferrule{C \subseteq D}{\clej{\Gamma}{\ev{elem}\,C\,D}{C}{D}} \\[6pt]
\rn{C-Union} & \inferrule{\clej{\Gamma}{f}{C_1}{D} \\ \clej{\Gamma}{g}{C_2}{D}}{\clej{\Gamma}{\ev{union}\,f\,g}{C_1 \cup C_2}{D}} \\[6pt]
\rn{C-Var} & \inferrule{\Gamma(x) = \capt{S}{C}}{\clej{\Gamma}{\ev{capvar}\,x}{\cset{x}}{C}} \\[10pt]
\rn{T-Var} & \inferrule{\Gamma(x) = T}{\tmj{\Gamma}{x}{T}} \\[6pt]
\rn{T-Val} & \inferrule{\valj{\Gamma}{v}{T}}{\tmj{\Gamma}{v}{T}} \\[6pt]
\rn{T-App} & \inferrule{\Gamma(x) = \capt{(\pit[z]{T}{T'})}{C} \\ \Gamma(y) = T}{\tmj{\Gamma}{x\,y}{T'[z \mapsto y]}} \\[6pt]
\rn{T-Let} & \inferrule{\tmj{\Gamma}{t}{T} \\ \tmj{\Gamma, x : T}{u}{\weak{T'}} \\ \clej{\Gamma, x : T}{f}{\uses{u}}{\weak{U}}}{\tmj{\Gamma}{\letc{x}{t}{u}{U}{f}}{T'}} \\[10pt]
\rn{V-Lam} & \inferrule{\tmj{\Gamma, x : T}{t}{T'} \\ \clej{\Gamma, x : T}{g}{\uses{t}}{\weak{A} \cup \cset{x}}}{\valj{\Gamma}{\lamc{A}{x}{T}{t}{g}}{\capt{(\pit[x]{T}{T'})}{A}}} \\[6pt]
\rn{V-Unit} & \inferrule{ }{\valj{\Gamma}{\unitv}{\Unit}}
\end{array}
\]
\color{black}
\caption{Typing of the fragment: capture evidence, terms and values. $\weak{X}$ is $X$ read under the new binder, which $X$ does not mention.}
\label{fig:mini-typing}
\end{figure}

Capture evidence is a proof term for an inclusion of capture sets. The judgment $\clej{\Gamma}{f}{C}{D}$ says that $f$ proves $C \le D$ in $\Gamma$. Figure~\ref{fig:mini-typing} gives its five rules, followed by the typing of terms and values. Rules are named by the conventions of Section~\ref{sec:intro}, and proofs cite them by name.

\rn{C-Refl} and \rn{C-Trans} make capture inclusion a preorder. \rn{C-Elem} is a syntactic inclusion, which a checker decides. \rn{C-Union} bounds a union by bounding each part. \rn{C-Var} is the documentation's rule that a capability is below the capture set of its type (Section~\ref{sec:cc-types}). No rule replaces a capture variable by another set, as \rn{C-Var} does for a term variable. A capture variable stands for itself.

\paragraph{Typing.}
The term rules are report~1's, restricted to variables, with capture sets and their evidence added. \rn{T-App} applies a variable to a variable, and the result type instantiates the codomain at the argument. The fragment has no subsumption, so \rn{T-App} needs the argument at the parameter's type exactly. In the full calculus report~1's casts supply subsumption (Section~\ref{sec:target-delta}). Two rules carry capture evidence of their own, in the brackets that erasure removes.

\rn{V-Lam} types a closure $\lamc{A}{x}{T}{t}{g}$ at $\capt{(\pit[x]{T}{T'})}{A}$, so its annotation becomes the capture set of its type. The closing evidence $g$ proves $\uses{t} \le \weak{A} \cup \cset{x}$. What the body uses is covered by the annotation and the parameter.

\rn{T-Let} types $\letc{x}{t}{u}{U}{f}$. Its avoidance evidence $f$ proves $\uses{u} \le \weak{U}$ under the binder of $x$, and the declared set $U$ does not mention $x$. This is avoidance (Section~\ref{sec:cc-uses}). The body may use $x$, but only through what $U$ covers. Why declare $U$ instead of computing it from $u$? The uses of $u$ may mention $x$, which is not in scope outside the let. Evidence turns them into a set that does not mention $x$, and the let records that set as data.

\paragraph{Use sets.}
A use set is what a term may use when it runs. Figure~\ref{fig:mini-syntax} computes it by structural recursion, because every set a binder declares is data. A value uses nothing, so creating a closure uses nothing. What the body of a closure uses is charged at each call, through the annotation of the closure (Section~\ref{sec:mini-step}). A let uses what its bound term uses and the set its body declares.

\paragraph{Inspection.}
A step inspects a variable when it reads the value stored there. In the fragment only a call does, and it inspects its function: $\insp{x\,y} = x$. An argument is used but not inspected. Effect safety (Theorem~\ref{thm:mini-effect}) speaks about the variables that a run inspects.

\begin{example}[Typing the example]\label{ex:mini-typing}
Each capability is typed with closing evidence that is a syntactic inclusion. The body $w$ of $v_{\mathit{log}}$ uses $\cset{w}$, and $\ev{elem}\,\cset{w}\,\cset{\kappa_1, w}$ proves $\cset{w} \le \cset{\kappa_1} \cup \cset{w}$. So $v_{\mathit{log}}$ has type $\CapT{\kappa_1}$, and in the same way $v_{\mathit{console}}$ has type $\CapT{\kappa_2}$.

The inner closure $v_{\mathit{inner}}$ is typed in $\Gamma_0, y : \Unit, r : \Unit$. Its body $\mathit{console}\,z$ has type $\Unit$ by \rn{T-App}. It uses $\cset{\mathit{console}, z}$, which $\ev{elem}$ puts below $\cset{\mathit{console}} \cup \cset{z}$. The body of $v_{\mathit{c1}}$ is a let. Its bound term $\mathit{log}\,y$ has type $\Unit$. Its body $v_{\mathit{inner}}$ is a value and uses nothing, so the let declares $\cempty$ with avoidance evidence $\ev{elem}\,\cempty\,\cempty$. The type of $v_{\mathit{inner}}$ does not mention $r$, so the let has that type. The let uses $\uses{\mathit{log}\,y} \cup \cempty = \cset{\mathit{log}, y}$, and the closing evidence of $v_{\mathit{c1}}$ puts that below $\cset{\mathit{log}} \cup \cset{y}$. Hence $\valj{\Gamma_0}{v_{\mathit{c1}}}{T_{\mathit{c1}}}$.

In $\Gamma_1$ the program $\mathit{prog}_1$ binds $\mathit{inner}$ at $\capt{(\pit[z]{\Unit}{\Unit})}{\cset{\mathit{console}}}$, the codomain of $T_{\mathit{c1}}$ at $\mathit{unit}$. Its body $\mathit{inner}\,\mathit{unit}$ uses $\cset{\mathit{inner}, \mathit{unit}}$, and its avoidance evidence is
\[
\begin{array}{@{}lcl@{}}
f_1 & = & \ev{union}\ (\ev{trans}\ (\ev{capvar}\,\mathit{inner})\ (\ev{elem}\,\cset{\mathit{console}}\,\cset{\mathit{console}, \mathit{log}})) \\
& & \phantom{\ev{union}}\ (\ev{trans}\ (\ev{capvar}\,\mathit{unit})\ (\ev{elem}\,\cempty\,\cset{\mathit{console}, \mathit{log}})) .
\end{array}
\]
It is built in steps. By \rn{C-Var}, $\ev{capvar}\,\mathit{inner}$ proves $\cset{\mathit{inner}} \le \cset{\mathit{console}}$, the capture set of the type of $\mathit{inner}$. By \rn{C-Elem} and \rn{C-Trans} the first half proves $\cset{\mathit{inner}} \le \cset{\mathit{console}, \mathit{log}}$. The type of $\mathit{unit}$ is $\Unit$, so $\ev{capvar}\,\mathit{unit}$ proves $\cset{\mathit{unit}} \le \cempty$, and the second half proves $\cset{\mathit{unit}} \le \cset{\mathit{console}, \mathit{log}}$. \rn{C-Union} joins the halves into $\cset{\mathit{inner}} \cup \cset{\mathit{unit}} \le \cset{\mathit{console}, \mathit{log}}$.

The program $\mathit{prog}_2$ drops $\mathit{inner}$. Its body $\mathit{unit}$ uses $\cset{\mathit{unit}}$, and $\ev{capvar}\,\mathit{unit}$ proves $\cset{\mathit{unit}} \le \cempty$. Both programs have type $\Unit$ in $\Gamma_1$, and
\[
\uses{\mathit{prog}_1} = \cset{\mathit{c1}, \mathit{unit}, \mathit{console}, \mathit{log}} ,
\qquad
\uses{\mathit{prog}_2} = \cset{\mathit{c1}, \mathit{unit}} .
\]
A declared set may be larger than the body needs. $\mathit{prog}_1$ declares $\mathit{log}$, which its body never calls, and evidence only has to show an inclusion. These typings and use sets are decided.
\end{example}

\FloatBarrier
\subsection{What capture evidence means}\label{sec:mini-canon}

Capture evidence gets no normal form. What a piece of capture evidence says is a fact about leaves, and that fact is all that the rest of the proof uses. The base case is a variable.

\begin{lemma}[Leaves of a variable]\label{lem:mini-var}
For every $\Gamma$ and $x$ with $\Gamma(x) = \capt{S}{C}$, $\capeqv{\Gamma}{\cset{x}}{C}$.
\end{lemma}

\begin{proof}
Let $\Gamma = \Gamma', x : \capt{S}{C}, \Gamma''$. By the first three equations for leaves, $\leaves{\Gamma}{\cset{x}} = \leaves{\Gamma'}{C}$. The set $C$ mentions only binders of $\Gamma'$, so $\leaves{\Gamma}{C} = \leaves{\Gamma'}{C}$.
\end{proof}

\begin{theorem}[Capture evidence includes leaves]\label{thm:mini-cap}
For every $\Gamma$, if $\clej{\Gamma}{f}{C}{D}$ then $\caple{\Gamma}{C}{D}$.
\end{theorem}

\begin{proof}[Proof sketch]
By induction on the derivation, one case per rule.
\begin{itemize}
\item \rn{C-Refl}: $\subseteq$ is reflexive.
\item \rn{C-Trans}: the two induction hypotheses compose, since $\subseteq$ is transitive.
\item \rn{C-Elem}: the leaves of a set are the union of the leaves of its atoms, and every atom of $C$ is an atom of $D$.
\item \rn{C-Union}: $\leaves{\Gamma}{C_1 \cup C_2} = \leaves{\Gamma}{C_1} \cup \leaves{\Gamma}{C_2}$ by the first two equations, and each part is covered by $D$ by an induction hypothesis.
\item \rn{C-Var}: Lemma~\ref{lem:mini-var}. \qedhere
\end{itemize}
\end{proof}

The fragment has no capture members, so the theorem holds in every context, and Section~\ref{sec:target-canon} shows where the full calculus needs the store.
